\documentclass[10pt]{article}
\usepackage[margin=20mm]{geometry}
\usepackage{amsmath,amssymb,fontspec,unicode-math,xurl,hyperref}
\IfFontExistsTF{DejaVu Serif}{\setmainfont{DejaVu Serif}}{\setmainfont{Latin Modern Roman}}
\setmathfont{latinmodern-math.otf}
\setlength{\parindent}{0pt}\setlength{\parskip}{6pt}\sloppy
\title{Technical Verification of MA-TSD}
\author{AI-assisted technical review}
\date{}
\begin{document}
\maketitle
Abstract: The Gaussian marginal construction and deterministic reverse-step algebra of MA-TSD pass the checks described here. Two aspects of the training rationale need a more precise interpretation: the hybrid-loss upper bound depends on whether sequence errors are summed or averaged, and the exact Gaussian variational term has a matrix-weighted error rather than the unweighted denoising error. A literal Euclidean reading of the hybrid proof admits a counterexample, whereas the mean-squared-error convention used in the released code admits a valid bound. These alternatives prevent a supported Major or Minor defect judgment in this review. The outcome is Needs Evidence within the declared scope; empirical results were not reproduced.\par
\section*{Paper identity and scope}
The reviewed paper is Chenxi Wang, Linxiao Yang, Zhixian Wang, Liang Sun, and Yi Wang, A Non-isotropic Time Series Diffusion Model with Moving Average Transitions, ICML 2025, PMLR 267, pp. 65144–65166, paper I25-43547 (local number 0178) [paper]. The source is the staged 23-page proceedings PDF. References to pages below mean one-based PDF pages, which coincide with the displayed article page numbers.\par
The claim-focused review covered the forward marginal and conditional Gaussian construction in Section 4.2 and Appendix B.3.2; the deterministic reverse algebra in Section 4.3 and Appendix B.3.3; the denoising objective in Equation (15); and the hybrid bound in Section 4.4 and Appendix B.3.4. Experiment descriptions were read for task definitions and the relation of the theory to the implementation. Relevant released source files were inspected at commit `6b80af723e5ac285f27ace2e63a284e694b90cc0` [code]. No author code, training, data loading, or benchmark evaluation was executed.\par
This is not a full review of every numerical result or a certification of the entire method. The coverage ledger records partially inspected and unexamined material separately.\par
\section*{Definitions and checks that passed}
The paper normalizes each nonconstant target series \(x_0\in\mathbb R^L\) using scalar temporal mean \(\mu\) and standard deviation \(\sigma\), giving \(z_0=(x_0-\mu)/\sigma\). Its marginal is \[q(z_t\mid z_0)=\mathcal N(K_tz_0,\beta_t^2 I).\] The interpolated nonoverlapping moving-average matrices \(K_t\) need not be invertible. For \(\beta_t>0\) and \(0\leq\eta_t^2\leq\beta_{t-1}^2\), define \[a_t=\frac{\sqrt{\beta_{t-1}^2-\eta_t^2}}{\beta_t}.\] Equation (14) is the Gaussian conditional with mean \(K_{t-1}z_0+a_t(z_t-K_tz_0)\) and covariance \(\eta_t^2 I\).\par
Conditional on \(z_0\), write \(z_t=K_tz_0+\beta_t\epsilon\) and add independent Gaussian innovation \(\eta_t\xi\). Then \[z_{t-1}=K_{t-1}z_0+
\sqrt{\beta_{t-1}^2-\eta_t^2}\,\epsilon+\eta_t\xi.\] Its mean and covariance are exactly \(K_{t-1}z_0\) and \(\beta_{t-1}^2 I\). This verifies the claimed marginal consistency without requiring an inverse for \(K_t\). The independent arithmetic check with \(\beta_{t-1}=0.6,\beta_t=0.8,\eta_t=0.3\) returns the required covariance coefficient \(0.36\). A subscript inconsistency in the displayed covariance expansion in Equation (21) does not invalidate this direct calculation.\par
For \(\eta_t=0\), dividing the reverse step by \(\beta_{t-1}\) yields the discrete identity in Equation (24): \[\frac{z_{t-1}}{\beta_{t-1}}-\frac{z_t}{\beta_t}
 =\left(\frac{K_{t-1}}{\beta_{t-1}}-\frac{K_t}{\beta_t}\right)
   f_\theta(z_t,t,c).\] This verifies the rearrangement. It does not by itself prove convergence of a chosen continuous-time discretization or accuracy after skipping steps.\par
The normalization argument in Appendix B.3.1 is valid when \(K_t\mathbf 1=\mathbf 1\): scalar shifts commute with such a matrix, and standard deviations are shift invariant. This check does not certify every claimed variance-preservation interpretation of the dataset-based noise schedule.\par
\section*{Training claims requiring clarification}
Problem 1 (The hybrid bound depends on the sequence-loss convention).  The literal Euclidean-norm proof and the mean-squared-error implementation lead to different constants. Assessment: Needs Evidence, with a conditional counterexample and a valid alternative reading.\par
\subsection*{Claim, source, and core dependency}
Section 4.4, page 6, states that appropriately chosen coefficients in \[L_{\rm hybrid}=\lambda_z L_z+\lambda_\mu L_\mu+\lambda_\sigma L_\sigma\] make this objective an upper bound on the conditional diffusion loss without instance normalization. Appendix B.3.4, page 15, Equations (26)–(38), supplies the proof and prescribes a common coefficient based on maxima of normalized target amplitude and predicted standard deviation.\par
This is a theoretical rationale for jointly training the normalized denoiser and the two statistic predictors, not an unrelated appendix observation. A demonstrated failure under the intended convention would concern that training guarantee. The issue cannot be graded merely from its location or the ease of changing a constant.\par
\subsection*{Literal Euclidean reading: exact counterexample}
As written, Equations (15) and (26) use \(\|\cdot\|_2^2\) for length-\(L\) vectors, while the errors in the scalar mean and standard deviation are scalar squares. Let the dataset consist of the single length-eight series \[x_0=z_0=(1,-1,1,-1,1,-1,1,-1)^\top.\] Its temporal population mean is zero and standard deviation is one. Let \(f_\theta(z_t,t,c)=z_0\) for every input, let \(g^\sigma_\omega(c)=1\), and let \(g^\mu_\omega(c)=1\). A constant-output denoiser is admissible, and the standard-deviation predictor satisfies the proposed cap at the dataset maximum \(\sigma_{\max}=1\). A singleton dataset and constant predictors also make the following calculation valid after expectation over all diffusion steps and noises.\par
The reconstructed series is \(\widehat x_0=z_0+\mathbf 1\). Thus \[L_x=\|\mathbf 1\|_2^2=8,\qquad L_z=0,\qquad
L_\sigma=0,\qquad L_\mu=1.\] Even using the generous squared-maxima interpretation of the proposed coefficient, \(\lambda=\max\|z_0\|_\infty^2+\max|g^\sigma_\omega|^2+1=3\), the hybrid loss is three. The purported inequality becomes \(8\leq3\).\par
Equation reference: eq:sumcorrect. The source of this discrepancy is the broadcast of scalar errors. With Euclidean vector errors, the exact decomposition is \[\widehat x_0-x_0=
\widehat\sigma(f_\theta-z_0)+z_0(\widehat\sigma-\sigma)
 +(\widehat\mu-\mu)\mathbf 1,\] so its norm obeys \[
\|\widehat x_0-x_0\|_2
\leq |\widehat\sigma|\,\|f_\theta-z_0\|_2
 +\|z_0\|_2|\widehat\sigma-\sigma|
 +\sqrt L|\widehat\mu-\mu|.\] The corresponding factors are absent from the substitutions leading to Equations (30)–(34). This calculation uses neither a constant series with undefined normalization nor an uncapped variance predictor.\par
\subsection*{Counterevidence: the released code uses averaged losses}
The paper describes its fitting objective as mean squared error. More decisively, the released `src/models/matsd.py` assigns `F.mse\_loss` to the reconstruction loss and also uses it for the scalar mean and standard-deviation losses in `\_get\_loss`. The temporal statistic predictors in `src/backbone/mlp.py` have output length one. The default mean reduction therefore gives an averaged sequence error, not the summed Euclidean error used in the literal calculation above.\par
Equation reference: eq:meancorrect. Under that convention, define \[\ell_{z,\mathrm{rms}}=L^{-1/2}\|f_\theta-z_0\|_2,\qquad
\ell_{x,\mathrm{rms}}=L^{-1/2}\|\widehat x_0-x_0\|_2,\] and let \(A\) and \(B\) bound \(\|z_0\|_\infty\) and \(|\widehat\sigma|\) on the domain being claimed. Dividing [eq:sumcorrect] (\#eq:sumcorrect) by \(\sqrt L\) and using \(\|z_0\|_2/\sqrt L\leq A\) gives \[\ell_{x,\mathrm{rms}}\leq
B\ell_{z,\mathrm{rms}}+A|\widehat\sigma-\sigma|
+|\widehat\mu-\mu|.\] Cauchy–Schwarz now proves \[
L_{x,\mathrm{MSE}}\leq
(B^2+A^2+1)
\bigl(L_{z,\mathrm{MSE}}+L_\sigma+L_\mu\bigr).\] For the example above, this reads \(1\leq3\) and is valid.\par
This is an admissible reading of the intended method, supported by actual source inspection. It prevents treating the literal-norm contradiction as a demonstrated defect of the implemented training objective. Conversely, the code’s average reduction does not make the written Euclidean calculation literally correct. An explicit convention in the mathematical statement is needed.\par
The bounded-predictor condition is part of [eq:meancorrect] (\#eq:meancorrect); it should not be silently inferred from an unconstrained linear network output. The reviewed backbone does not visibly impose the cap contemplated in Appendix B.3.4. The paper’s experiments also use coefficients equal to one, whereas the theorem discussion conditions the bound on suitable coefficients. We do not infer that the paper promised its bound for those experimental settings.\par
\subsection*{Impact and precisely verified alternatives}
The supported result is conditional: the displayed summed-norm reading fails; the averaged-MSE reading has a valid pointwise and expected bound when the stated amplitude caps hold. For summed errors, Cauchy–Schwarz applied to [eq:sumcorrect] (\#eq:sumcorrect) instead gives the sufficient common coefficient \(B^2+LA^2+L\). Either statement is precise. Neither verifies author adoption or the historical implementation used for the reported results.\par
A clarification of the intended losses and boundedness domain would resolve this concern. No benchmark deterioration, invalid Gaussian sampler, or universal failure of hybrid training is established.\par
Problem 2 (An exact variational term is matrix weighted).  Equation (15) is a defensible denoising surrogate, but it is not in general an algebraic scalar reweighting of the Gaussian KL term induced by Equation (14). Assessment: Needs Evidence about the intended strength of the derivation, with the distinction proved below.\par
\subsection*{Original statement and exact calculation}
Equation reference: eq:kl. Section 4.3, page 5, introduces the ordinary squared-error objective in Equation (15) by analogy with DDIM. The preceding background, pages 2–3, motivates training through Gaussian variational terms. For \(\eta_t>0\), substitute \(f_\theta\) for \(z_0\) in Equation (14), as the paper prescribes. At fixed \((z_t,z_0,c)\), the two Gaussian covariances coincide and their means differ by \[A_t(z_0-f_\theta),\qquad A_t=K_{t-1}-a_tK_t.\] The exact conditional divergence is therefore \[
D_{\mathrm{KL}}\!\left(q(z_{t-1}\mid z_t,z_0)
 \,\|\,p_\theta(z_{t-1}\mid z_t)\right)
=\frac{\|A_t(z_0-f_\theta)\|_2^2}{2\eta_t^2}.\] The matrix \(A_t^\top A_t\) need not be a scalar multiple of the identity.\par
For an admissible explicit example, use length four, the two nonoverlapping pair averages followed by linear interpolation, and the final all-series average: \[K_1=\begin{pmatrix}
1/2&1/2&0&0\\3/8&3/8&1/8&1/8\\
1/8&1/8&3/8&3/8\\0&0&1/2&1/2
\end{pmatrix},
\qquad K_2=\tfrac14\mathbf1\mathbf1^\top.\] These matrices preserve constants and the temporal mean. For the normalized singleton series \((1,1,-1,-1)^\top\), the paper’s schedule gives \(\gamma_1=\sqrt{5/8}\), \(\beta_1^2=3/8\), and \(\beta_2^2=1\). Choose \(\eta_2^2=\beta_1^2\), which satisfies the admissible noise constraint. Then \(a_2=0\) and \(A_2=K_1\). The errors \[d_1=(1,-1,0,0)^\top,\qquad d_2=(1,1,0,0)^\top\] both have squared norm two, but [eq:kl] (\#eq:kl) equals zero for \(d_1\) and \(13/6\) for \(d_2\). Thus no scalar time weighting makes the individual KL term identical to the ordinary squared error in this example. This is a finite calculation supporting an exact algebraic distinction, not an author-code execution.\par
\subsection*{Alternative reading and limits}
The source says it derives a similar optimization target by analogy; it does not provide a detailed equality proof for the non-isotropic metric. An ordinary squared-error loss can instead be chosen as a surrogate. More directly, for a fixed deterministic positive-semidefinite matrix \(H=A_t^\top A_t\), the conditional mean \(m=\mathbb E[z_0\mid z_t,t,c]\) minimizes both ordinary MSE and \[\mathbb E[(z_0-f)^\top H(z_0-f)\mid z_t,t,c].\] Indeed, writing \(z_0-f=(z_0-m)+(m-f)\) eliminates the conditional cross term and leaves a nonnegative quadratic in \(m-f\). If \(H\) is singular, the weighted loss can have additional minimizers. With a restricted network and shared parameters, equivalence of optimizers is not automatic.\par
Consequently, [eq:kl] (\#eq:kl) does not refute the use of Equation (15), the marginal consistency proof, or the reported experiments. It rules out only an exact scalar-reweighting interpretation without additional assumptions. The report requests explicit identification of the surrogate objective and the level at which the analogy is intended, rather than assigning a Major issue from an interpretation the paper may not assert.\par
\section*{Overall assessment}
The checked Gaussian construction is valid under its variance and nonzero-denominator conditions. No supported core sampler failure or empirical error was established. The two concerns concern the formal interpretation of the training rationale and remain Needs Evidence because an admissible interpretation is supported by the text or implementation.\par
The staged manifest includes the appendix and a pinned repository snapshot. Potential additional OpenReview attachments were unresolved during source preparation. Later revisions and errata were not independently searched; none is claimed to have been excluded or incorporated. Determining historical experimental impact would require provenance and, if the actual objective changes, fresh training and evaluation.\par
Checked object | Inspection judgment | Repair / clarification | Remaining state\par
Gaussian marginal | No Issue Found in scope | Direct algebra verified | No empirical certification\par
Hybrid bound | Needs Evidence | Mean-MSE bound verified under caps | Norm convention/adoption unknown\par
KL versus MSE | Scope clarification; Needs Evidence | Exact metric derived; surrogate remains valid | Intended claim unresolved\par
Scientific decision | Pending | Separate from these checks | Human reviews empty\par
9 C. Wang, L. Yang, Z. Wang, L. Sun, and Y. Wang, A Non-isotropic Time Series Diffusion Model with Moving Average Transitions, PMLR 267 (2025), 65144–65166. Proceedings PDF, 23 pages. https://proceedings.mlr.press/v267/wang25dy.html. C. Wang et al., Moving-Average-Diffusion, released source commit `6b80af723e5ac285f27ace2e63a284e694b90cc0` (2025-06-23). https://github.com/WillWang1113/Moving-Average-Diffusion.\par
\section*{Supporting checks}
The reviewed original is `papers/0178/0178.pdf`, SHA-256 c81f5b2e2cd9d6921eb6b6bc1d06a223047435d4d141c0013916bfab4b62414e. It matches the staged source manifest. The existing `pdftotext -layout` utility extracted the paper.\par
The independent check, run with Python 3.13.5 and NumPy 2.4.3, imports no author code. Its output records the Euclidean sum eight, averaged MSE one, hybrid bound three, Gaussian covariance coefficient \(0.36\), and KL values zero and \(13/6\) for equal Euclidean errors. Training uses a training-only data loader to estimate its schedule in the inspected script; this static observation is not a full data-leakage audit.\par
No human review, author contact or author adoption is asserted.\par
\end{document}
